\documentclass[11pt,a4paper]{article}
\usepackage[T1]{fontenc}
\usepackage[utf8]{inputenc}
\usepackage{lmodern}
\usepackage[margin=2.5cm]{geometry}
\usepackage{amsmath,amssymb}
\usepackage{graphicx}
\usepackage{booktabs}
\usepackage{tabularx}
\usepackage{algorithm}
\usepackage{algpseudocode}
\usepackage[numbers,sort&compress]{natbib}
\usepackage[hidelinks]{hyperref}
\usepackage[capitalize]{cleveref}
\usepackage{caption}
\captionsetup{font=small,labelfont=bf}

%%%%%%%%% TITLE - PLEASE UPDATE TITLE HERE
\title{My Title}

\author{Asaf Markuza}
\date{}

\begin{document}

\maketitle

\begin{abstract}
%%%%%%%%% ABSTRACT %%%%%%%%%
Abstract here.
%%%%%%%%% ABSTRACT %%%%%%%%%
\end{abstract}

\section{Introduction}
\label{sec:intro}
%%%%%%%%% INTRODUCTION %%%%%%%%%
Physiological systems do not operate in isolation. The brain, heart, lungs, eyes and skeletal muscles continuously coordinate their activity, and the pattern of these interactions changes with physiological state. Network physiology studies this coordination by representing organ systems as nodes and their dynamical couplings as links of a time-varying network~\cite{Bashan2012NetworkPhysiology,Bartsch2015NetworkPhysiology,Ivanov2021NewField}. The approach has shown that the network of organ interactions reorganises abruptly across sleep stages~\cite{Bashan2012NetworkPhysiology,Bartsch2015NetworkPhysiology}, that cardio-respiratory coupling undergoes transition-like changes between states~\cite{Bartsch2012PhaseTransitions}, that brain-wave interactions evolve across physiological states~\cite{Liu2015PlasticityBrain,Liu2015MajorComponent}, that the brain--muscle network is stratified by sleep stage~\cite{Rizzo2020NetworkPhysiology,Rizzo2022DynamicNetworks}, and that brain rhythms and heart rate interact with characteristic time delays~\cite{Lin2016DelaycorrelationLandscape}.

Most of these results rest on one pairwise coupling measure, Time Delay Stability (TDS)~\cite{Bashan2012NetworkPhysiology,Bartsch2015NetworkPhysiology}. TDS slides a window over two signals, finds in each window the time delay $\tau_0$ that maximises their cross-correlation, and declares the pair coupled for as long as $\tau_0$ stays nearly constant over consecutive windows. A window is labelled stable when at least 4 of 5 neighbouring delays agree within a tolerance, and the TDS score is the percentage of stable windows. The method is simple, interpretable and directly tied to a physical quantity, the delay, which is why it has become the standard tool of the field.

TDS, however, depends on three parameters counted in samples: the window length $L$, the lag range $\tau_{\max}$ searched, and the tolerance $\mathrm{tol}$. The published values, $L = 60$, $\tau_{\max} = \pm 30$ and $\mathrm{tol} = \pm 1$, were chosen for physiological series resampled to 1\,Hz. Because they are counted in samples, their meaning changes with both the time scale of the system and the sampling rate: the same physical coupling recorded at a sixteen-fold higher rate is seen by a 60-sample window through a sixteenth of the time span. For a new system there is no rule to choose them, and different implementations expose and set them differently, which makes TDS results hard to compare and reproduce~\cite{Steinbrinker2024TdspyOpensource}. A fixed choice can also fail silently. We show below that when the true delay lies outside a fixed $\pm30$ lag range, TDS still returns high scores (72--96\,\%) while reporting a wrong delay at the edge of the search, often with the wrong sign: the TDS score measures how stable a delay is, not whether it is right.

The statistical root of the problem is well understood. The reliability of a correlation computed from autocorrelated data depends on the number of effectively independent samples, not on the raw sample count~\cite{Bartlett1946TheoreticalSpecification,Bayley1946EffectiveNumber,Pyper1998ComparisonMethods,Afyouni2019EffectiveDegrees}, and the accuracy of cross-correlation delay estimates depends on observation time and spectral content~\cite{Knapp1976GeneralizedCorrelation,Carter1987CoherenceTime}. A window of fixed length therefore carries less and less information the slower the signals are. Likewise, whether a given TDS score could arise by chance depends on each signal's own temporal structure, which is exactly what surrogate-data methods are designed to preserve while destroying coupling~\cite{Theiler1992TestingNonlinearity,Schreiber2000SurrogateTime,Lancaster2018SurrogateData}.

We build on these principles to propose \emph{adaptive Time Delay Stability} (aTDS), a set of three rules that set all TDS parameters from the signals themselves, without coupling labels or known delays:
\begin{enumerate}
\item \textbf{Window.} $L = 30 \cdot \max_i B_i$, where $B_i$ is the Bartlett variance-inflation factor of signal $i$'s autocorrelation~\cite{Bartlett1946TheoreticalSpecification}, so that each window holds about 30 independent samples of the slowest signal (at least 61 samples, at most 10\,\% of the record; optionally at least $2 d_{\max} + 1$ for a known longest plausible delay $d_{\max}$).
\item \textbf{Lag range.} $\tau_{\max} = \lfloor (L-1)/2 \rfloor$, the published ratio and the aliasing limit of the circular cross-correlation, capped at half the shortest detected rhythm period.
\item \textbf{Tolerance.} The widest tolerance whose chance score on circularly rotated, coupling-free copies of the signals, mean plus two standard errors, stays within a 5\,\% false-stability budget.
\end{enumerate}
The 4-of-5 stability rule of the original method is unchanged.

Our contributions are: (i) the three rules above, each with an explicit statistical rationale; (ii) a controlled evaluation on signals with known coupling and delay that isolates each rule and compares the full method with the published setting, a setting rescaled by hand with the true time-scale factor, and an oracle tuned on labelled data, using AUC~\cite{Hanley1982MeaningUse} with bootstrap intervals~\cite{Efron1979BootstrapMethods}; (iii) tests on long-memory~\cite{Davies1987TestsHurst,Peng1995QuantificationScaling}, rhythmic, bursting~\cite{Rulkov2002ModelingSpikingbursting}, mixed-speed and non-stationary signals and on one real sleep-EEG pair; and (iv) a per-link rotation test with false-discovery-rate control~\cite{Benjamini1995ControllingFalse} for networks. We scope our claims to these tests: aTDS automates a time-scale-consistent choice of TDS parameters with a false-alarm budget; it does not find optimal parameters, and its effect on sleep-stage contrasts in real multi-subject cohorts remains to be established.
%%%%%%%%% INTRODUCTION %%%%%%%%%

\section{Related Work}
\label{sec:related_work}
%%%%%%%%% RELATED WORK %%%%%%%%%
\paragraph{Delay-based coupling measures in physiology.}
TDS was introduced to quantify how pairs of physiological signals lock onto a stable time delay and to build networks of organ interactions across sleep stages~\cite{Bashan2012NetworkPhysiology,Bartsch2015NetworkPhysiology}. Subsequent work used it to characterise brain-wave networks~\cite{Liu2015PlasticityBrain}, the dominant components of dynamic organ networks~\cite{Liu2015MajorComponent}, and cortico--muscular interactions, where TDS was computed in 60\,s windows moved in 30\,s steps~\cite{Rizzo2020NetworkPhysiology,Rizzo2022DynamicNetworks}. The delay-correlation landscape generalises the idea by examining the full profile of correlation against delay for brain--heart interactions~\cite{Lin2016DelaycorrelationLandscape}, and cardio-respiratory coupling has been shown to change abruptly between sleep stages~\cite{Bartsch2012PhaseTransitions}. Across these studies the window, lag range and tolerance are fixed in samples or seconds and carried over from the original setting; the open-source TDSpy implementation exposes them as user parameters and notes that differing choices hamper reproducibility~\cite{Steinbrinker2024TdspyOpensource}. All of these methods inherit the classical theory of time-delay estimation by cross-correlation, whose accuracy depends on coherence and on the length of the observation window~\cite{Knapp1976GeneralizedCorrelation,Carter1987CoherenceTime}. Other coupling measures address different questions: Granger causality~\cite{Granger1969InvestigatingCausal} and transfer entropy~\cite{Schreiber2000MeasuringInformation} quantify directed predictive influence, convergent cross mapping~\cite{Sugihara2012DetectingCausality} targets coupling in nonlinear deterministic systems, and coupling-function methods reconstruct the form of interactions between oscillators~\cite{Stankovski2017CouplingFunctions}. These methods have their own embedding, lag and window parameters; aTDS is complementary to them and concerns the parameter choice of TDS itself.

\paragraph{Effective sample size of autocorrelated series.}
Bartlett showed that the variance of a correlation estimated between two autocorrelated series is inflated by a factor determined by the products of their autocorrelation functions~\cite{Bartlett1946TheoreticalSpecification}, and Bayley and Hammersley derived the corresponding effective number of independent observations of an autocorrelated series~\cite{Bayley1946EffectiveNumber}. These corrections are routinely used to adjust correlation tests in ecology~\cite{Pyper1998ComparisonMethods} and neuroimaging~\cite{Afyouni2019EffectiveDegrees}. In the TDS setting, the same reasoning implies that a fixed-length window carries fewer independent samples the slower the signals, so that the cross-correlation peak, and thus $\tau_0$, becomes unreliable. aTDS turns this observation into a design rule: it sets the window so that it always contains a fixed number of independent samples, using the Bartlett factor of each signal's autocorrelation.

\paragraph{Surrogate data and multiple testing.}
Surrogate data are the standard way to ask whether a statistic could have arisen without the property of interest~\cite{Theiler1992TestingNonlinearity,Schreiber2000SurrogateTime,Lancaster2018SurrogateData}. Phase-randomised and amplitude-adjusted surrogates preserve the spectrum and distribution of a series~\cite{Theiler1992TestingNonlinearity,Schreiber1996ImprovedSurrogate}; multivariate surrogates also preserve cross-spectra~\cite{Prichard1994GeneratingSurrogate}; and bivariate surrogates are needed to test for interdependence between two signals while keeping each signal's own structure~\cite{Andrzejak2003BivariateSurrogate}. The original TDS work builds its null by pairing signals across subjects within the same sleep stage~\cite{Bashan2012NetworkPhysiology}, which requires a cohort. aTDS uses a within-recording null, circular rotation of one signal relative to the other, which keeps both autocorrelations exactly and applies to a single system; it uses this null twice, to calibrate the tolerance under a fixed false-stability budget and to obtain a per-link $p$-value. Because a network of $n$ signals involves $n(n-1)/2$ simultaneous tests, the per-link $p$-values are combined with the Benjamini--Hochberg procedure, which controls the expected proportion of false links among those reported~\cite{Benjamini1995ControllingFalse}.
%%%%%%%%% RELATED WORK %%%%%%%%%

\section{Methods}
\label{sec:method}
%%%%%%%%% METHOD %%%%%%%%%
Methodology here.
%%%%%%%%% METHOD %%%%%%%%%

\section{Results}
\label{sec:experiments}
%%%%%%%%% EXPERIMENTS %%%%%%%%%
Experimental setup, results, and analysis here.
%%%%%%%%% EXPERIMENTS %%%%%%%%%

\section{Discussion}
\label{sec:discussion}
%%%%%%%%% DISCUSSION %%%%%%%%%
Discussion, limitations and future work here.
%%%%%%%%% DISCUSSION %%%%%%%%%

\section{Conclusion}
\label{sec:conclusion}
%%%%%%%%% CONCLUSION %%%%%%%%%
Conclusion here.
%%%%%%%%% CONCLUSION %%%%%%%%%

\section*{Code and data availability}
pyTDS, the aTDS experiment scripts, their raw results (JSON) and the figure scripts are available at \url{https://github.com/Marku-a/pyTDS}. All random seeds are fixed.

%%%%%%%%% REFERENCES %%%%%%%%%
\bibliographystyle{unsrtnat}
\bibliography{references}
%%%%%%%%% REFERENCES %%%%%%%%%

%%%%%%%%% Appendix %%%%%%%%%
% If you want to add appendix, uncomment the following lines:
% \appendix
% \section{xxx}
%%%%%%%%% Appendix %%%%%%%%%

\end{document}